Let $N=9$ and $\lambda$ be the partition $(3,2,1)$ and $\mu=(3,1,0)$. Let $a=(t^2,t)$ such that $a^+=a\cup \{0,0,0,0,\infty,\infty,\infty\}$, in the notation of this paper: \href{https://arxiv.org/abs/0902.1321v3}{arXiv:0902.1321v3}. Let $T\in \mathrm{SYT}(\lambda/\mu;a)$ be the tableau with filling of the skew shape with $t$ in the bottom row and $t^2$ in the cell above and to the right of this cell. In the language of the just mentioned paper the tableau is a diagonally increasing skew tableau.

According to Theorem~4.2 in the same paper there exists a linear series $x$, with certain properties. Let $f$ be the map $\mathbb{P}^1\rightarrow \mathbb{P}^2$ corresponding to $x$. Consider $(f,\{t^2,t\})$ as an element in the moduli space of stable maps, with marked points at $t^2$, $t$.

In the limit $t\rightarrow 0$ in the moduli space of stable maps this becomes a degree $4$ plane curve with a degree zero component $X$ attached at $c_\infty$ that has one marked point on it at $c_t$ and one at $c_0$, corresponding to the marked points $t$, $t^2$, respectively. There is a degree $1$ component attached to $X$ at $c_*$.

Assume that $X$ is parametrized such that $c_\infty$ is at $c=\infty$, $c_0$ is at $c=0$, and $c_t$ is at $c=1$. At what value of $c$ is the node $c_*$?
